\section{Robótica, cronograma de la AGI e impacto económico: optimismo tecnológico junto con fricción real}

\subsection{El mundo de altas restricciones de la IA encarnada}

El diálogo sobre robotics resume su juicio como «hay oportunidad, pero el umbral es extremadamente alto». En los sistemas digitales, un fallo puede revertirse con rapidez; en los sistemas físicos, un error se convierte en un incidente de seguridad. Lex resumió esta restricción en una frase: \textit{``embodied systems are almost allowed to fail never''}. Esto hace que el despliegue de robots avance con un ritmo claramente más lento que el de los agentes puramente de software.

\begin{importantbox}{Por qué la conducción autónoma y la automatización industrial se despliegan antes}
\begin{itemize}[leftmargin=1.5em]
    \item Los límites de la tarea son más claros y las métricas de evaluación pueden definirse;
    \item el entorno de operación es relativamente controlable y la incertidumbre puede reducirse mediante infraestructura;
    \item el retorno económico es cuantificable y la ruta entre inversión y beneficio es más definida.
\end{itemize}
\end{importantbox}

\subsection{Cronograma de la AGI: la raíz de la divergencia está en la definición}

La entrevista no ofrece un cronograma unificado; la causa principal es que la definición de AGI no es consistente. Si se define como «completar la mayoría de las tareas de la economía digital», el plazo podría ser más cercano; si se define como «superar de manera estable al ser humano en todas las dimensiones cognitivas», el plazo es claramente más lejano. Nathan insiste repetidamente en la característica de la «jagged intelligence»: el modelo es extremadamente fuerte en algunas tareas y muy frágil en otras.

\begin{warningbox}{La etiqueta de «programador sobrehumano» puede ocultar brechas sistémicas}
Equiparar unas cuantas puntuaciones altas en benchmarks con «poder sustituir un rol de ingeniería complejo» produce un juicio erróneo. La producción real de software incluye colaboración organizacional, negociación de requisitos, comunicación entre equipos y responsabilidad de mantenimiento a largo plazo; nada de esto puede cubrirse con la generación de código en una sola ronda.
\end{warningbox}

\subsection{Resumen de la sección}

En 2026, la discusión sobre la AGI debe mantener una postura de ingeniería: primero observar el cierre de tareas concretas, y solo después hablar de etiquetas grandilocuentes. La robótica y la automatización integral no llegarán de forma lineal, sino que se irán infiltrando escenario por escenario.
